% Generated draft: requires mathematical and source review.
\documentclass[11pt]{article}
\usepackage[margin=25mm]{geometry}
\usepackage{fontspec}
\setmainfont{Latin Modern Roman}
\usepackage{amsmath,amssymb,mathtools,unicode-math}
\setmathfont{Latin Modern Math}
\usepackage{enumitem,needspace,xurl,hyperref}
\setlength{\parindent}{0pt}
\setlength{\parskip}{5pt}
\setlength{\emergencystretch}{3em}
\newcommand{\N}{\mathbb{N}}
\newcommand{\Z}{\mathbb{Z}}
\newcommand{\Q}{\mathbb{Q}}
\newcommand{\R}{\mathbb{R}}
\newcommand{\C}{\mathbb{C}}
\newcommand{\F}{\mathbb{F}}
\newcommand{\A}{\mathbb{A}}
\newcommand{\PP}{\mathbb{P}}
\newcommand{\E}{\mathbb{E}}
\newcommand{\eps}{\varepsilon}
\newcommand{\dd}{\,\mathrm{d}}
\begin{document}
\section*{UnsolvedMath ID: 30002430}
\textbf{Source problem code:} catalogue entry
\subsection{Definitions and mathematical statement}
Write \(F_1,\ldots,F_R\) for the homogeneous forms, \(d\) for their common degree, and \(n\) for their number of variables. The supplied record asks:
\begin{quote}
For every nonsingular system of \(R\) homogeneous forms of equal degree, is
\[
B=R-1,
\]
where \(B\) denotes the dimension of its singular locus?
\end{quote}
The record does not specify the coefficient field, the permitted values of \(n,d,R\), the meaning of nonsingularity, or which singular locus and dimension convention define \(B\). These omissions prevent an unambiguous mathematical formulation.

For clarity, two potentially different loci can be distinguished without selecting either as the intended definition. If the forms lie in \(K[x_1,\ldots,x_n]\) for a field \(K\), let \(\overline K\) be an algebraic closure and write
\[
J_F(x)=\left(\frac{\partial F_i}{\partial x_j}(x)\right)_{\substack{1\leq i\leq R\\1\leq j\leq n}}.
\]
One possible intended locus is the affine Jacobian rank locus
\[
S_{\mathrm{aff}}(F)=\{x\in\overline K^{\,n}:\operatorname{rank}J_F(x)<R\},
\]
which imposes no equations \(F_i(x)=0\). Another possible meaning concerns the singular locus of the common projective zero scheme
\[
X_F=\{[x]\in\mathbb P^{n-1}_{\overline K}:F_1(x)=\cdots=F_R(x)=0\}.
\]
These definitions are supplied here only to expose the unresolved interpretation. If nonsingularity means that \(X_F\) is smooth, its own singular locus is empty; that reading does not identify \(B\) with the dimension of \(S_{\mathrm{aff}}(F)\). If the intended convention is \(B=\dim S_{\mathrm{aff}}(F)\), the question becomes whether this dimension is always \(R-1\) for systems satisfying the intended nonsingularity hypothesis. Neither this convention nor the necessary field, degree, dimension, and nonsingularity hypotheses are fixed by the supplied record. Affine versus projective dimension and the convention for an empty locus also remain unresolved.
\subsection{Short English statement}
The supplied question asks whether the singular locus associated with every nonsingular system of \(R\) equal-degree homogeneous forms has dimension exactly \(R-1\). A precise statement requires clarification of the field, nonsingularity hypothesis, and intended singular locus. The supplied dated curation labels the problem open; this draft does not independently verify that status.
\subsection{Sources}
\begin{itemize}
\item Supplied local UnsolvedMath registry record, permanent ID \(30002430\), exact source problem code \texttt{catalogue entry}, titled \emph{Singular Locus Dimension of Nonsingular Form Systems}; collection: Oberwolfach Reports Open Problems. No individual proposer is identified. The proposed year is recorded as 2013.
\item Supplied report citation: \emph{Analytic Number Theory (2014)}, Oberwolfach Reports, pp.~2963--3037, DOI \texttt{10.4171/owr/2013/51}. Source URL and supplied literature-source link, labeled ``original OWR report'': \url{https://doi.org/10.4171/owr/2013/51}. The supplied background describes these reports as records of mathematical workshops and their open-problem discussions. The report itself was not consulted for this draft.
\item Literature assessment dated exactly \texttt{2026-08-21}: the supplied curation labels the problem open, reports that a known argument gives \(B\leq R-1\), and says that equality \(B=R-1\) is not established in the generality of the extracted question. This is supplied dated curation, not an independently verified literature assessment. Its stated source is the original OWR report linked above; the record expressly requests independent verification of time-sensitive status and cited claims.
\item The registry labels the statement \texttt{corrected\_verified} and describes its editing as normalization of notation, typography, and wording without changing mathematical content. This supplied label does not resolve the formulation ambiguities identified here or constitute mathematical review of this draft.
\end{itemize}
\end{document}
